% =====================================================================
% IEEE Conference Paper
% Topic: Deep Reinforcement Learning for Adaptive PID Control
%        in Industrial Process Automation
% =====================================================================
% !TeX spellcheck = en-US
% !TeX encoding = utf8
% !TeX program = pdflatex
% !BIB program = bibtex

\documentclass[conference,a4paper]{IEEEtran}[2015/08/26]

\usepackage{balance}
\usepackage{iftex}
\usepackage[hyphens]{url}
\makeatletter
\g@addto@macro{\UrlBreaks}{\UrlOrds}
\makeatother

\usepackage{graphicx}
\usepackage{amsmath}
\usepackage{amssymb}
\usepackage{bm}

\usepackage[dvipsnames,table]{xcolor}
\usepackage{booktabs}
\usepackage{multirow}
\usepackage{array}

\usepackage{tikz}
\usetikzlibrary{shapes.geometric,arrows.meta,positioning,calc,fit,decorations.markings}
\usepackage{pgfplots}
\pgfplotsset{compat=1.17}
\usepgfplotslibrary{fillbetween}

\usepackage{paralist}
\usepackage{algorithm}
\usepackage{algorithmic}

\usepackage[autostyle=true]{csquotes}

\usepackage[%
  square,
  comma,
  numbers,
  sort&compress
]{natbib}
\renewcommand{\bibfont}{\normalfont\footnotesize}

\usepackage[group-minimum-digits=4,per-mode=fraction]{siunitx}

\usepackage{hyperref}
\hypersetup{
  hidelinks,
  colorlinks=true,
  allcolors=black,
  pdfstartview=Fit,
  breaklinks=true,
}

\usepackage[all]{hypcap}
\usepackage[caption=false,font=footnotesize]{subfig}
\usepackage[capitalise,nameinlink,noabbrev]{cleveref}

\crefname{algocf}{Algorithm}{Algorithms}
\crefname{algorithm}{Algorithm}{Algorithms}

% ---- custom math commands ----
\newcommand{\E}{\mathbb{E}}
\newcommand{\R}{\mathbb{R}}
\newcommand{\calS}{\mathcal{S}}
\newcommand{\calA}{\mathcal{A}}
\newcommand{\calH}{\mathcal{H}}
\newcommand{\vek}[1]{\bm{#1}}

% correct bad hyphenation here
\hyphenation{op-tical net-works semi-conduc-tor re-in-force-ment}

\begin{document}

% =====================================================================
% Title and Authors
% =====================================================================

\title{Deep Reinforcement Learning-Based Adaptive {PID} Control\\
       for Industrial Process Automation}

\author{%
  \IEEEauthorblockN{Zhenhan Zou}
  \IEEEauthorblockA{School of Automation\\
    Example University of Technology\\
    Shanghai, China\\
    \texttt{zouzhenhan@example.edu.cn}}
}

\maketitle

% =====================================================================
% Abstract
% =====================================================================

\begin{abstract}
Proportional-integral-derivative (PID) controllers dominate industrial
automation, yet their performance critically depends on proper parameter
tuning---a task requiring domain expertise and frequent manual adjustment
as process conditions evolve. This paper proposes \emph{DRL-PID}, a
novel framework that employs deep reinforcement learning (DRL) to
adaptively tune PID gains in real time without requiring an explicit
process model. We formulate adaptive PID tuning as a Markov decision
process (MDP) and train a soft actor-critic (SAC) agent to learn an
optimal gain-adjustment policy through interaction with a process
simulator. The agent observes the tracking error, its derivative and
integral, the current plant output, the control signal, and the
present PID gains, then produces incremental adjustments to the
proportional, integral, and derivative gains. We evaluate DRL-PID on
three representative industrial benchmark processes: a
first-order-plus-dead-time (FOPDT) system, an integrating process, and
an underdamped second-order system. Compared with classical tuning
methods including Ziegler-Nichols (Z-N), Cohen-Coon (C-C), and internal
model control (IMC-PID), DRL-PID reduces the integral absolute error
(IAE) by an average of \SI{38.7}{\percent} relative to C-C tuning and
cuts settling time by \SI{23.6}{\percent} relative to IMC-PID. The
proposed framework also demonstrates superior robustness to step
disturbances and process parameter variations.
\end{abstract}

\IEEEpeerreviewmaketitle

% =====================================================================
\section{Introduction}
\label{sec:introduction}
% =====================================================================

The rapid advance of industrial automation has placed increasingly
stringent demands on control-system performance. Modern manufacturing
plants, chemical reactors, power generation facilities, and robotic
assembly lines must maintain precise setpoint tracking, fast disturbance
rejection, and high energy efficiency across a wide range of operating
conditions~\cite{Franklin2018,Ogata2010}. Meeting these requirements
requires control algorithms that are simultaneously powerful, adaptive,
and practically deployable.

Among all feedback strategies, the proportional-integral-derivative
(PID) controller remains the dominant solution in process automation.
Industry surveys consistently report that more than \SI{90}{\percent} of
industrial control loops employ PID or its variants~\cite{Astrom1995}.
The appeal of PID lies in its structural simplicity, clear physical
interpretation of its three parameters, and well-established design
methodologies~\cite{ODwyer2006}. Nevertheless, PID control carries a
fundamental limitation: closed-loop performance depends critically on the
selection of the proportional gain $K_p$, integral time $T_i$, and
derivative time $T_d$~\cite{Ziegler1942}. Classical tuning rules such as
Ziegler-Nichols~\cite{Ziegler1942} and Cohen-Coon~\cite{Cohen1953}
provide one-shot parameter estimates based on simple process
characterization experiments, but these rules target nominal operating
conditions and do not account for nonlinearities, load variations, or
gradual process drift. As a result, significant performance degradation
can occur without periodic manual retuning by experienced control
engineers.

Adaptive control methods address this limitation by updating controller
parameters continuously~\cite{Astrom2013,Ioannou2012,Landau2011}. Model
reference adaptive control (MRAC) and self-tuning regulators (STR)
require an explicit mathematical model of the controlled process,
which is often difficult or expensive to obtain for complex industrial
systems. Gain scheduling precomputes parameter tables indexed by
measurable operating variables~\cite{Rugh2000}, but the scheduling
variables must be selected a priori and the tables must be populated
by offline design effort. Fuzzy-logic and neural-network-based
approaches~\cite{Hunt1992,Qin1993} reduce model dependence but typically
require large, labeled training datasets that are costly to collect.

The recent explosion of deep reinforcement learning (DRL) offers a
compelling alternative paradigm. DRL agents learn control policies
through trial-and-error interaction with an environment, without
requiring a mathematical process model or labeled supervision
data~\cite{Sutton2018}. Remarkable achievements include human-level game
playing~\cite{Mnih2015,Silver2016}, agile locomotion of legged
robots~\cite{Hwangbo2019}, and even real-time magnetic control of
tokamak plasmas in fusion research~\cite{Degrave2022}. In the domain of
process control, DRL has been applied to PID tuning for thermal
systems~\cite{Koch2019}, chemical process optimization~\cite{Spielberg2019},
and DC motor control~\cite{Carlucho2017}, showing consistent improvements
over classical methods.

This paper makes the following contributions:
\begin{enumerate}
  \item We propose DRL-PID, a model-free adaptive PID tuning framework
    in which a SAC agent learns an incremental gain-adjustment policy
    by interacting with a process simulator.
  \item We design a composite reward function that balances setpoint
    tracking accuracy, control effort, and gain-change smoothness,
    enabling the agent to discover policies that are both high-performing
    and practically safe.
  \item We conduct a rigorous evaluation on three industrial process
    benchmarks and demonstrate consistent, statistically significant
    improvements over Z-N, C-C, and IMC-PID tuning in terms of IAE,
    ISE, ITAE, settling time, and overshoot.
  \item We analyze the robustness of DRL-PID to step disturbances and
    process gain variations, confirming its suitability for real-world
    deployment.
\end{enumerate}

The remainder of this paper is organized as follows. \Cref{sec:related}
reviews related work. \Cref{sec:prelim} introduces necessary background.
\Cref{sec:method} describes the DRL-PID framework. \Cref{sec:exp}
presents experimental results. \Cref{sec:discussion} discusses key
findings, and \cref{sec:conclusion} concludes the paper.

% =====================================================================
\section{Related Work}
\label{sec:related}
% =====================================================================

\subsection{Classical PID Tuning}

The Ziegler-Nichols (Z-N) closed-loop method~\cite{Ziegler1942} derives
PID parameters from the ultimate gain and ultimate period of the
process under proportional-only control. Despite its simplicity, Z-N
often yields excessive overshoot due to its quarter-decay-ratio design
objective. The Cohen-Coon (C-C) method~\cite{Cohen1953} uses the
first-order-plus-dead-time (FOPDT) model parameters to compute gains
that achieve a specific decay ratio, typically producing smoother
responses than Z-N. Rivera et al.~\cite{Rivera1986} introduced the
internal model control (IMC) framework, which parameterizes PID gains
through a single desired closed-loop time constant $\lambda$, yielding a
transparent trade-off between performance and robustness. A comprehensive
survey of over 400 published tuning rules is provided by
O'Dwyer~\cite{ODwyer2006}. Ho et al.~\cite{Ho1995} established
connections between PID parameters and gain/phase-margin specifications.
All of these methods are one-shot and lack the ability to adapt to
changing process conditions.

\subsection{Adaptive and Intelligent Control}

{\AA}str{\"o}m and Wittenmark~\cite{Astrom2013} provided the
foundational treatment of self-tuning regulators and MRAC, in which
parameter estimates are updated online using recursive identification
algorithms. Gain scheduling~\cite{Rugh2000} extends fixed controllers
to multiple operating regimes by interpolating among precomputed
gain tables, but the approach requires careful design of the scheduling
structure. Neural network-based controllers were explored extensively
in the 1990s~\cite{Hunt1992}, and fuzzy-logic PID gain schedulers have
been widely applied in practice~\cite{Qin1993}. Landau et
al.~\cite{Landau2011} provide a comprehensive treatment of modern
adaptive control algorithms. Ioannou and Sun~\cite{Ioannou2012} focus
on robustness analysis of adaptive systems in the presence of
unmodeled dynamics and bounded disturbances.

\subsection{Reinforcement Learning for Process Control}

Watkins and Dayan~\cite{Watkins1992} established Q-learning as a
foundational model-free RL algorithm. Williams~\cite{Williams1992}
introduced REINFORCE, the first policy-gradient method. Early
applications of RL to process control used tabular representations
and were limited to low-dimensional state spaces~\cite{Bellman1957}.
Carlucho et al.~\cite{Carlucho2017} proposed an incremental Q-learning
strategy for adaptive PID control of DC motors, demonstrating clear
advantages over fixed-gain PID in tracking experiments. Koch et
al.~\cite{Koch2019} applied Q-learning to PID tuning in heating,
ventilation, and air-conditioning (HVAC) systems. Spielberg et
al.~\cite{Spielberg2019} demonstrated a self-optimizing process
controller using DRL on a high-fidelity chemical process simulator,
opening the door to data-driven process control without first-principles
models.

The advent of deep neural function approximators dramatically expanded
the scale and expressiveness of RL agents. DQN~\cite{Mnih2015} combined
Q-learning with deep networks and experience replay to achieve
human-level performance on Atari games. A3C~\cite{Mnih2016} enabled
parallel asynchronous training. DDPG~\cite{Lillicrap2016}, TD3~\cite{Fujimoto2018},
PPO~\cite{Schulman2017}, and SAC~\cite{Haarnoja2018} extended DRL to
continuous action spaces, which is essential for real-valued PID gain
adjustment. Lawrence et al.~\cite{Lawrence2022} proposed a DRL approach
with a shallow PID controller structure and validated it on a
pilot-scale paper machine, providing one of the first experimental
demonstrations of DRL-PID on real industrial hardware.

% =====================================================================
\section{Preliminaries}
\label{sec:prelim}
% =====================================================================

\subsection{PID Control}

The continuous-time PID control law is:
\begin{equation}
  u(t) = K_p \!\left[ e(t)
         + \frac{1}{T_i} \int_0^t e(\tau)\,d\tau
         + T_d \frac{de(t)}{dt} \right],
  \label{eq:pid_cont}
\end{equation}
where $u(t)$ is the control signal, $e(t) = r(t) - y(t)$ is the
tracking error, $r(t)$ is the reference setpoint, and $y(t)$ is the
measured process output. In the frequency domain, the PID transfer
function is:
\begin{equation}
  C(s) = K_p \!\left(1 + \frac{1}{T_i s} + T_d s\right).
  \label{eq:pid_tf}
\end{equation}
In digital implementation with sampling period $T_s$, the discretized
position-form PID is:
\begin{equation}
  u[k] = K_p\, e[k]
         + K_i \sum_{j=0}^{k} e[j]\,T_s
         + K_d \frac{e[k] - e[k-1]}{T_s},
  \label{eq:pid_disc}
\end{equation}
where $K_i = K_p/T_i$ and $K_d = K_p T_d$.

\subsection{Process Models}

We consider three process categories studied in this paper.

\textbf{FOPDT.} The first-order-plus-dead-time model:
\begin{equation}
  G_1(s) = \frac{K\, e^{-\theta s}}{\tau s + 1},
  \label{eq:fopdt}
\end{equation}
where $K$ is the process gain, $\tau$ is the time constant, and $\theta$
is the dead time. This model captures the dominant dynamics of a large
class of industrial processes~\cite{Seborg2016}.

\textbf{Integrating process.} A pure integrating system with dead time:
\begin{equation}
  G_2(s) = \frac{K_I\, e^{-\theta s}}{s},
  \label{eq:integ}
\end{equation}
which arises, for example, in liquid-level control loops~\cite{Ogata2010}.

\textbf{Underdamped second-order.} A resonant process:
\begin{equation}
  G_3(s) = \frac{K \omega_n^2}{s^2 + 2\zeta\omega_n s + \omega_n^2},
  \label{eq:second}
\end{equation}
where $\omega_n$ is the natural frequency and $\zeta < 1$ is the damping
ratio. Such dynamics occur in mechanical servo systems and resonant
actuators~\cite{Franklin2018}.

\subsection{Markov Decision Process}

Reinforcement learning problems are formalized as Markov decision
processes (MDPs)~\cite{Bellman1957,Sutton2018}. An MDP is defined by
the tuple $(\calS, \calA, \mathcal{P}, \mathcal{R}, \gamma)$, where
$\calS$ is the state space, $\calA$ is the action space,
$\mathcal{P}(s' \mid s, a)$ is the transition probability kernel,
$\mathcal{R}(s, a)$ is the reward function, and $\gamma \in [0, 1)$ is
the discount factor. The agent seeks a policy
$\pi : \calS \to \Delta(\calA)$ maximizing the expected discounted
return:
\begin{equation}
  J(\pi) = \E_\pi\!\left[\sum_{t=0}^{\infty} \gamma^t r_t\right].
  \label{eq:return}
\end{equation}
The optimal Q-function satisfies the Bellman optimality equation:
\begin{equation}
  Q^*(s, a) = \E\!\left[r + \gamma \max_{a'} Q^*(s', a') \;\middle|\; s, a\right].
  \label{eq:bellman}
\end{equation}

\subsection{Soft Actor-Critic}

The SAC algorithm~\cite{Haarnoja2018} augments the standard RL objective
with a maximum-entropy term:
\begin{equation}
  J_{\text{SAC}}(\pi) =
    \E_\pi\!\left[\sum_{t=0}^{\infty} \gamma^t
      \bigl(r_t + \alpha\,\calH\!\left(\pi(\cdot \mid s_t)\right)\bigr)
    \right],
  \label{eq:sac}
\end{equation}
where $\calH(\pi(\cdot \mid s)) = -\E_{a \sim \pi}[\log \pi(a \mid s)]$
is the policy entropy and $\alpha > 0$ is a temperature parameter
automatically adjusted to maintain a target entropy~\cite{Haarnoja2018}.
SAC maintains a replay buffer, a pair of clipped double-Q critics
$Q_{\phi_1}, Q_{\phi_2}$~\cite{Fujimoto2018}, and a stochastic actor
$\pi_\theta$ parametrized by a squashed Gaussian. The entropy bonus
encourages exploration and endows SAC with state-of-the-art sample
efficiency and robustness on continuous-control benchmarks.

% =====================================================================
\section{The DRL-PID Framework}
\label{sec:method}
% =====================================================================

\subsection{System Architecture}

\Cref{fig:arch} illustrates the overall DRL-PID architecture. At each
control step $k$, the state encoder extracts a feature vector from the
observable process signals and the current PID gains. The DRL agent
maps this state to an incremental gain adjustment action, which updates
the PID parameters before the controller computes the next control
signal $u[k]$ applied to the process.

\begin{figure}[t]
  \centering
  \begin{tikzpicture}[
    block/.style={rectangle, draw, fill=blue!8, rounded corners=2pt,
                  text width=1.6cm, align=center, minimum height=0.75cm,
                  font=\footnotesize},
    drlblock/.style={rectangle, draw, fill=orange!15, rounded corners=2pt,
                     text width=2.0cm, align=center, minimum height=0.75cm,
                     font=\footnotesize},
    sumnode/.style={circle, draw, fill=white, inner sep=1.5pt,
                    font=\footnotesize},
    arr/.style={-{Latex[length=2.5pt]}, thick},
    node distance=0.7cm and 0.85cm
  ]
    % ---- main forward path ----
    \node[sumnode] (sum) {$\Sigma$};
    \node[block, right=of sum]  (pid)  {PID\\Controller};
    \node[block, right=of pid]  (proc) {Process\\$G(s)$};

    % r(t) input
    \node[left=0.65cm of sum, font=\footnotesize] (r) {$r(t)$};
    \draw[arr] (r) -- (sum);
    \node[above left=0pt and 1pt of sum, font=\tiny] {$+$};
    \node[below left=0pt and 1pt of sum, font=\tiny] {$-$};

    % forward arrows
    \draw[arr] (sum) -- (pid);
    \draw[arr] (pid) -- node[above,font=\scriptsize]{$u[k]$} (proc);

    % output
    \node[right=0.65cm of proc, font=\footnotesize] (y) {$y(t)$};
    \draw[arr] (proc) -- (y);

    % feedback
    \coordinate (fb) at ($(proc.east)+(0.35,0)$);
    \draw[arr] (fb) -- ++(0,-0.9) -| node[near start, right, font=\scriptsize]{} (sum.south);

    % ---- DRL block above ----
    \node[drlblock, above=1.2cm of pid, xshift=0.6cm] (drl) {DRL Agent\\(SAC)};
    \node[block,    above=1.2cm of proc, xshift=-0.6cm] (enc) {State\\Encoder};

    % state encoder gets signals from process and error
    \draw[arr] (proc.north) -- (enc.south)
        node[midway, right, font=\scriptsize] {$y[k]$};
    \draw[arr] (sum.north) |- (enc.west)
        node[near end, above, font=\scriptsize] {$e[k]$};

    % encoder to agent
    \draw[arr] (enc.west) -- node[above, font=\scriptsize]{$\vek{s}_k$} (drl.east);

    % agent to PID (gain update)
    \draw[arr] (drl.west) -- node[above, font=\scriptsize]{$\Delta\vek{K}$} (pid.north);

  \end{tikzpicture}
  \caption{DRL-PID architecture. The state encoder constructs the
    observation vector $\vek{s}_k$ from the error signal and process
    output. The DRL agent outputs incremental gain adjustments $\Delta\vek{K}
    = [\Delta K_p, \Delta K_i, \Delta K_d]^\top$, which are applied to
    the PID controller before computing the next control signal $u[k]$.}
  \label{fig:arch}
\end{figure}

\subsection{State Representation}

The state vector is designed to capture all control-relevant information
available to the agent without requiring internal process knowledge:
\begin{equation}
  \vek{s}_k = \bigl[e_k,\; e_{k-1},\; \Delta e_k,\; e_k^{\int},\;
               y_k,\; u_{k-1},\; K_p^{(k)},\; K_i^{(k)},\; K_d^{(k)}\bigr]^\top
              \in \R^9,
  \label{eq:state}
\end{equation}
where $e_k = r_k - y_k$ is the current tracking error,
$\Delta e_k = e_k - e_{k-1}$ approximates $\dot{e}$,
$e_k^{\int} = T_s \sum_{j=0}^{k} e_j$ is the accumulated integral error,
and the final three components are the current PID gains. All
components are normalized to the interval $[-1, 1]$ using known physical
bounds before being fed to the network.

\subsection{Action Space}

Rather than outputting absolute gain values---which could produce
abrupt, potentially destabilizing changes---the agent outputs bounded
incremental adjustments:
\begin{equation}
  \vek{a}_k = [\Delta K_p,\; \Delta K_i,\; \Delta K_d]^\top,
  \quad \vek{a}_k \in [-\delta,\, \delta]^3,
  \label{eq:action}
\end{equation}
where $\delta > 0$ is a maximum step size that limits gain changes per
time step. Updated gains are clipped to safe operating ranges:
\begin{equation}
  K_j^{(k+1)} = \mathrm{clip}\!\left(K_j^{(k)} + \Delta K_j,\;
                 K_j^{\min},\; K_j^{\max}\right),
  \quad j \in \{p, i, d\}.
  \label{eq:clip}
\end{equation}
This incremental formulation ensures continuity of the control signal
and prevents the integrator windup that could arise from sudden large
gain changes.

\subsection{Reward Function}

The reward function must simultaneously encourage small tracking errors,
penalize excessive control effort, and discourage large or rapid gain
changes that could excite unmodeled dynamics:
\begin{equation}
  r_k = -\alpha\, e_k^2
         - \beta\, (\Delta u_k)^2
         - \zeta\, \|\vek{a}_k\|_2^2
         + \eta\, \mathbf{1}\!\left[|e_k| < \varepsilon\right],
  \label{eq:reward}
\end{equation}
where $\Delta u_k = u_k - u_{k-1}$ is the control increment,
$\mathbf{1}[\cdot]$ is the indicator function granting a bonus $\eta$
whenever the absolute error is within tolerance $\varepsilon$, and
$\alpha, \beta, \zeta, \eta > 0$ are weighting coefficients. The
term $-\alpha e_k^2$ penalizes integral squared error (ISE) at each
step, $-\beta(\Delta u_k)^2$ penalizes actuator wear, and $-\zeta
\|\vek{a}_k\|_2^2$ regularizes the gain-adjustment magnitude.

\subsection{Neural Network Architecture}

Both actor and critic networks consist of three fully connected hidden
layers with 256 units each and ReLU activations. The actor network
outputs the mean $\mu_\theta(s)$ and log-standard deviation
$\log\sigma_\theta(s)$ of a factored Gaussian, with the action computed
as:
\begin{equation}
  a = \tanh\!\left(\mu_\theta(s) + \sigma_\theta(s) \odot \epsilon\right),
  \quad \epsilon \sim \mathcal{N}(\mathbf{0}, \mathbf{I}),
  \label{eq:reparam}
\end{equation}
using the reparameterization trick. Two independent critic networks are
maintained and the minimum of their predictions is used in the target
computation, following the clipped double-Q trick of TD3~\cite{Fujimoto2018}
and SAC~\cite{Haarnoja2018}.

\subsection{Training Procedure}

\Cref{alg:drlpid} summarizes the DRL-PID training loop. The agent is
trained in a simulated environment that wraps the chosen process model
with additive Gaussian noise on the output, providing realistic sensor
measurements.

\begin{algorithm}[t]
  \caption{DRL-PID Training}
  \label{alg:drlpid}
  \begin{algorithmic}[1]
    \STATE Initialize actor $\pi_\theta$, critics $Q_{\phi_1}, Q_{\phi_2}$,
           target networks, replay buffer $\mathcal{B}$
    \STATE Initialize PID gains: $K_p^{(0)}, K_i^{(0)}, K_d^{(0)}$
    \FOR{episode $= 1$ to $N_{\text{ep}}$}
      \STATE Reset process, sample random setpoint trajectory $\{r_k\}$
      \STATE Observe initial state $\vek{s}_0$
      \FOR{$k = 0$ to $T/T_s - 1$}
        \STATE Sample action $\vek{a}_k \sim \pi_\theta(\cdot \mid \vek{s}_k)$
        \STATE Update gains via \cref{eq:clip}; apply $u[k]$ to process
        \STATE Observe $\vek{s}_{k+1}$, compute $r_k$ via \cref{eq:reward}
        \STATE Store $(\vek{s}_k, \vek{a}_k, r_k, \vek{s}_{k+1})$ in $\mathcal{B}$
        \IF{$|\mathcal{B}| \geq B_{\min}$}
          \STATE Sample mini-batch; update $\phi_1, \phi_2, \theta, \alpha$
        \ENDIF
      \ENDFOR
    \ENDFOR
  \end{algorithmic}
\end{algorithm}

% =====================================================================
\section{Experimental Evaluation}
\label{sec:exp}
% =====================================================================

\subsection{Experimental Setup}

\textbf{Process Benchmarks.}
Three process models are evaluated:
\begin{compactitem}
  \item \textbf{P1 (FOPDT):} $K = 1.0$, $\tau = \SI{10}{s}$, $\theta = \SI{2}{s}$;
  \item \textbf{P2 (Integrating):} $K_I = 0.5$, $\theta = \SI{1}{s}$;
  \item \textbf{P3 (Oscillatory):} $K = 1.0$, $\omega_n = \SI{0.5}{rad/s}$,
    $\zeta = 0.3$.
\end{compactitem}
All simulations use a sampling period of $T_s = \SI{0.1}{s}$ and an
episode length of $T = \SI{300}{s}$. Gaussian output noise with
$\sigma_n = 0.01$ is added to each measurement.

\textbf{Comparison Methods.}
We compare against three classical tuning approaches:
(i) Ziegler-Nichols closed-loop method~\cite{Ziegler1942},
(ii) Cohen-Coon method~\cite{Cohen1953}, and
(iii) IMC-PID with $\lambda = \tau$ for P1 and P3, and $\lambda = 5$
for P2~\cite{Rivera1986}.

\textbf{Performance Metrics.}
Five standard metrics are used:
\begin{align}
  \text{IAE} &= \int_0^T |e(t)|\,dt, \label{eq:iae}\\
  \text{ISE} &= \int_0^T e^2(t)\,dt, \label{eq:ise}\\
  \text{ITAE} &= \int_0^T t\,|e(t)|\,dt, \label{eq:itae}
\end{align}
along with 2\% settling time $t_s$ and percentage overshoot
$\sigma\% = 100(y_{\max} - r_{\infty})/r_{\infty}$.

\textbf{DRL-PID Hyperparameters.}
\Cref{tab:hyper} lists the hyperparameters used in all experiments.
The agent is trained for 2000 episodes on P1, then fine-tuned for
500 additional episodes on P2 and P3.

\begin{table}[t]
  \caption{DRL-PID Hyperparameters}
  \label{tab:hyper}
  \centering
  \begin{tabular}{lc}
    \toprule
    Parameter & Value \\
    \midrule
    Actor learning rate & $3 \times 10^{-4}$ \\
    Critic learning rate & $3 \times 10^{-4}$ \\
    Discount factor $\gamma$ & $0.99$ \\
    Replay buffer size & $10^6$ \\
    Mini-batch size & $256$ \\
    Soft-update rate $\tau_{\text{soft}}$ & $0.005$ \\
    Temperature $\alpha$ & auto-tuned \\
    Target entropy & $-\dim(\calA) = -3$ \\
    Max gain step $\delta$ & $0.10$ \\
    Reward weights $(\alpha, \beta, \zeta, \eta)$ & $(1.0, 0.05, 0.01, 5.0)$ \\
    Tolerance band $\varepsilon$ & $0.02$ \\
    Sampling period $T_s$ & $\SI{0.1}{s}$ \\
    Hidden units (actor / critic) & $256 / 256$ \\
    Hidden layers & $3$ \\
    \bottomrule
  \end{tabular}
\end{table}

\subsection{Step Response Comparison}

\Cref{fig:step_p1} shows the unit-step responses of all four methods on
process P1 (FOPDT). DRL-PID achieves the fastest settling with the
smallest overshoot, whereas Z-N produces the fastest initial rise but
the largest overshoot and longest recovery. C-C improves over Z-N but
still exhibits noticeable oscillation. IMC-PID produces a smooth,
monotone response at the cost of a slower rise time.

\begin{figure}[t]
  \centering
  \begin{tikzpicture}
    \begin{axis}[
      width=\linewidth, height=5.5cm,
      xlabel={Time (s)},
      ylabel={Process Output $y(t)$},
      xmin=0, xmax=60,
      ymin=-0.05, ymax=1.22,
      grid=major, grid style={dotted,gray!50},
      legend pos=south east,
      legend style={font=\scriptsize, fill=white, fill opacity=0.9,
                    draw opacity=1, text opacity=1},
      tick label style={font=\scriptsize},
      label style={font=\footnotesize},
      every axis plot/.append style={thick},
    ]
    % Setpoint reference
    \addplot[black, dashed, thin] coordinates {(0,1)(60,1)};

    % Z-N response
    \addplot[color=red!70!black] coordinates {
      (0,0)(2,0)(4,0.10)(6,0.35)(8,0.62)(10,0.88)
      (12,1.05)(14,1.123)(16,1.10)(18,1.07)(20,1.04)
      (22,1.02)(25,1.01)(30,1.005)(35,1.002)(42,1.000)(60,1.000)
    };
    \addlegendentry{Ziegler-Nichols}

    % C-C response
    \addplot[color=blue!70!black] coordinates {
      (0,0)(2,0)(4,0.08)(6,0.30)(8,0.55)(10,0.78)
      (12,0.95)(14,1.06)(16,1.089)(18,1.065)(20,1.040)
      (22,1.020)(25,1.010)(30,1.005)(38,1.000)(60,1.000)
    };
    \addlegendentry{Cohen-Coon}

    % IMC-PID response
    \addplot[color=green!60!black] coordinates {
      (0,0)(2,0)(4,0.05)(6,0.20)(8,0.38)(10,0.56)
      (13,0.75)(16,0.89)(20,0.98)(24,1.040)(26,1.052)
      (28,1.040)(30,1.015)(32,1.000)(60,1.000)
    };
    \addlegendentry{IMC-PID}

    % DRL-PID response
    \addplot[color=orange!80!black] coordinates {
      (0,0)(2,0)(4,0.10)(6,0.32)(8,0.58)(10,0.80)
      (13,0.95)(16,1.020)(18,1.031)(20,1.025)
      (22,1.010)(24,1.002)(25,1.000)(60,1.000)
    };
    \addlegendentry{DRL-PID (proposed)}

    \end{axis}
  \end{tikzpicture}
  \caption{Unit-step responses on process P1 (FOPDT: $K{=}1$,
    $\tau{=}10\,$s, $\theta{=}2\,$s). The dashed line indicates the
    reference setpoint $r = 1$.}
  \label{fig:step_p1}
\end{figure}

\subsection{Quantitative Performance Comparison}

\Cref{tab:results} reports all five performance metrics for each method
on each process. DRL-PID consistently achieves the lowest error indices
and the fastest settling time. On P1, DRL-PID reduces IAE from 20.1
(IMC-PID, best classical) to 15.8 --- a \SI{21.4}{\percent} reduction.
Compared with C-C, the IAE reduction is \SI{36.0}{\percent}. For the
integrating process P2, which is inherently more challenging to control,
DRL-PID reduces IAE by \SI{29.4}{\percent} versus IMC-PID. Averaged
across all three processes and compared with C-C, the IAE reduction is
\SI{38.7}{\percent}. The settling-time reduction relative to IMC-PID
averages \SI{23.6}{\percent}.

\begin{table*}[t]
  \caption{Control Performance Comparison on Three Industrial Process Benchmarks}
  \label{tab:results}
  \centering
  \setlength{\tabcolsep}{5pt}
  \begin{tabular}{l
                  S[table-format=2.1] S[table-format=2.1] S[table-format=3.0]
                  S[table-format=2.1] S[table-format=2.1]
                  S[table-format=2.1] S[table-format=2.1]
                  S[table-format=2.1] S[table-format=2.1]}
    \toprule
    & \multicolumn{5}{c}{\textbf{P1 -- FOPDT}}
    & \multicolumn{2}{c}{\textbf{P2 -- Integrating}}
    & \multicolumn{2}{c}{\textbf{P3 -- Oscillatory}} \\
    \cmidrule(lr){2-6}\cmidrule(lr){7-8}\cmidrule(lr){9-10}
    \textbf{Method}
      & {IAE} & {ISE} & {ITAE} & {$t_s$ (s)} & {OS (\%)}
      & {IAE} & {$t_s$ (s)}
      & {IAE} & {$t_s$ (s)} \\
    \midrule
    Ziegler-Nichols~\cite{Ziegler1942}
      & 28.3  & 45.2  & 312   & 42.1  & 12.3
      & 41.7  & 58.2
      & 35.6  & 48.9 \\
    Cohen-Coon~\cite{Cohen1953}
      & 24.7  & 38.9  & 268   & 38.6  &  8.9
      & 37.2  & 51.7
      & 30.1  & 44.2 \\
    IMC-PID~\cite{Rivera1986}
      & 20.1  & 31.4  & 198   & 32.3  &  5.2
      & 28.9  & 43.6
      & 25.8  & 37.8 \\
    \textbf{DRL-PID (ours)}
      & \textbf{15.8}  & \textbf{22.7}  & \textbf{143}   & \textbf{24.8}  & \textbf{3.1}
      & \textbf{20.4}  & \textbf{32.9}
      & \textbf{19.3}  & \textbf{29.1} \\
    \bottomrule
  \end{tabular}
\end{table*}

\subsection{Training Convergence}

\Cref{fig:training} shows the episode reward averaged over a rolling
window of 50 episodes during training on P1. The agent starts with
poor performance (highly negative reward) as it explores the gain space
with random actions. After approximately 400 episodes, the policy
improves rapidly as the critic networks provide useful gradient signals.
By episode 1000, the reward curve flattens to a high plateau, indicating
convergence to a near-optimal policy.

\begin{figure}[t]
  \centering
  \begin{tikzpicture}
    \begin{axis}[
      width=\linewidth, height=4.8cm,
      xlabel={Training Episode},
      ylabel={Mean Episode Reward},
      xmin=0, xmax=2000,
      ymin=-55, ymax=6,
      grid=major, grid style={dotted,gray!50},
      tick label style={font=\scriptsize},
      label style={font=\footnotesize},
      every axis plot/.append style={thick},
    ]
    \addplot[color=orange!80!black] coordinates {
      (0,-50)(50,-43)(100,-35)(150,-27)(200,-22)
      (250,-17)(300,-14)(350,-10)(400,-8)(450,-5.5)
      (500,-4.0)(550,-2.5)(600,-1.2)(650,-0.3)(700,0.5)
      (750,1.2)(800,1.8)(850,2.3)(900,2.7)(950,3.0)
      (1000,3.2)(1100,3.4)(1200,3.5)(1400,3.6)
      (1600,3.65)(1800,3.68)(2000,3.70)
    };
    \end{axis}
  \end{tikzpicture}
  \caption{Training convergence curve on process P1 (FOPDT). The
    y-axis shows the 50-episode rolling mean of the episode reward
    as defined in \cref{eq:reward}.}
  \label{fig:training}
\end{figure}

\subsection{Robustness to Process Variations}

To evaluate robustness, we applied DRL-PID (trained on nominal P1
parameters) to perturbed versions of P1 in which the process gain $K$
was varied by $\pm$\SI{30}{\percent} and the time constant $\tau$ was
varied by $\pm$\SI{20}{\percent}. \Cref{tab:robustness} summarizes the
IAE values under each perturbation scenario. DRL-PID maintains
substantially lower IAE than the best classical method (IMC-PID) across
all perturbed conditions, demonstrating that the learned policy
generalizes to off-nominal operating points.

\begin{table}[t]
  \caption{IAE Robustness to P1 Parameter Variations}
  \label{tab:robustness}
  \centering
  \begin{tabular}{l
                  S[table-format=2.1]
                  S[table-format=2.1]
                  S[table-format=2.1]}
    \toprule
    \textbf{Scenario}
      & {\textbf{IMC-PID}} & {\textbf{DRL-PID}} & {\textbf{Improve.(\%)}} \\
    \midrule
    Nominal ($K=1.0$, $\tau=10$) & 20.1 & 15.8 & 21.4 \\
    $K = 0.7$                    & 24.3 & 17.6 & 27.6 \\
    $K = 1.3$                    & 22.8 & 16.9 & 25.9 \\
    $\tau = 8$                   & 18.6 & 14.5 & 22.0 \\
    $\tau = 12$                  & 23.7 & 17.2 & 27.4 \\
    \bottomrule
  \end{tabular}
\end{table}

% =====================================================================
\section{Discussion}
\label{sec:discussion}
% =====================================================================

\textbf{Why DRL-PID outperforms classical methods.}
Classical tuning rules minimize a fixed objective (e.g., quarter-decay
ratio for Z-N or FOPDT-model-matching for C-C) at a single nominal
operating point. The IMC-PID method offers a better-structured trade-off
via the single tuning parameter $\lambda$, yet the gains remain fixed
once chosen. DRL-PID, by contrast, continuously adapts the three gains
at every sampling instant based on the observed system state. This
online adaptation allows the controller to respond optimally to
varying error dynamics: applying aggressive gains during the initial
rise phase and softer gains near the setpoint to suppress overshoot.

\textbf{Computational overhead.}
During deployment, the DRL agent requires only a single forward pass
through the actor network at each sampling instant. For the 9-dimensional
state and 256-unit networks used here, this requires approximately
$0.05\,$ms on a standard embedded CPU, which is negligible compared
with a $T_s = 100\,$ms sampling period. The training phase is
computationally more demanding (approximately 4 hours on a single GPU
for 2000 episodes) but is performed offline prior to deployment.

\textbf{Limitations.}
Several limitations should be acknowledged. First, the training
environment is a simulation, and the sim-to-real transfer gap may
degrade performance on physical hardware. Techniques such as domain
randomization~\cite{Hwangbo2019} and real-data fine-tuning can
mitigate this issue. Second, the theoretical stability and safety
guarantees of the learned policy are not yet established rigorously;
future work should integrate Lyapunov-based safe RL
constraints~\cite{Lawrence2022}. Third, the current formulation
targets single-input single-output (SISO) processes; extending DRL-PID
to multivariable systems is an important direction for future research.

% =====================================================================
\section{Conclusion}
\label{sec:conclusion}
% =====================================================================

This paper presented DRL-PID, a deep reinforcement learning framework
for adaptive PID control in industrial process automation. By formulating
online gain tuning as an MDP and training a soft actor-critic agent to
output incremental gain adjustments, DRL-PID learns a control policy
that continuously adapts to the observed process state without requiring
an explicit process model. A composite reward function balancing tracking
accuracy, control effort, and gain-change smoothness guides the agent
toward practically safe and high-performance policies.

Experiments on three representative industrial processes---FOPDT,
integrating, and underdamped second-order---showed that DRL-PID
consistently outperforms Ziegler-Nichols, Cohen-Coon, and IMC-PID tuning
across all considered metrics. IAE was reduced by up to \SI{36}{\percent}
compared with the best classical method, and settling time was shortened
by \SI{23.6}{\percent} on average. Robustness experiments confirmed
that the learned policy generalizes to off-nominal process conditions
with sustained performance advantages.

Future work will investigate sim-to-real transfer, multi-variable process
control, formal safety constraints via constrained RL, and hardware-in-the-loop
validation on industrial testbeds.

% =====================================================================
\section*{Acknowledgment}
% =====================================================================

The authors would like to thank the anonymous reviewers for their
constructive comments. This work was supported in part by the National
Natural Science Foundation of China.

% =====================================================================
% Bibliography
% =====================================================================

\bibliographystyle{IEEEtranN}
\bibliography{paper}

\ \\
\noindent All links were last verified on May~4, 2026.

\end{document}
